% META: {"id": "TNSI-ALGO-EVAL-B", "chapitre": "TNSI-ALGORITHMIQUE", "type_objet": "evaluation", "version": "B", "capacites": ["T-ALGO-02A", "T-ALGO-02B", "T-ALGO-02C", "T-ALGO-02D", "T-ALGO-03", "T-ALGO-04", "T-ALGO-05"], "mode_creation": "ex_nihilo", "sources_inspiration": [], "duree_min": 45, "status": "generated"}
\section*{Évaluation B --- Graphes et diviser pour régner}
\textit{Durée : 45 minutes. Ordinateur avec interpréteur Python autorisé.}

\baremeIndicatif{Ex. 1 : 7 pts (parcours, cycle, chemin) ; Ex. 2 : 6 pts (diviser pour régner) ; Ex. 3 : 4 pts (programmation dynamique) ; Ex. 4 : 3 pts (Boyer-Moore)}

\bigskip

\textbf{Exercice 1} --- \textit{Parcours, cycle et chemin} \hfill (7 points)

On donne le graphe non orienté :
\begin{python}
graphe = {
    "P": ["Q", "R"],
    "Q": ["P", "S"],
    "R": ["P", "S"],
    "S": ["Q", "R"],
}
\end{python}
\begin{enumerate}
  \item (2 pts) Donner l'ordre de visite des sommets par un parcours en largeur (BFS) depuis \texttt{"P"}, en supposant que les voisins sont explorés dans l'ordre où ils apparaissent dans la liste.
  \item (2 pts) Donner l'ordre de visite par un parcours en \textbf{profondeur} (DFS) depuis
  \texttt{"P"}, avec la meme convention, et dire quelle structure de donnees distingue les
  deux parcours.
  \item (2 pts) Ce graphe contient-il un cycle ? Justifier en citant un cycle si oui.
  \item (1 pt) Ecrire une fonction \lstinline{existe_chemin(graphe, depart, arrivee)} qui dit
  s'il existe un chemin entre deux sommets, et donner un chemin de \texttt{"P"} a
  \texttt{"S"}.
\end{enumerate}

% BEGIN-VERIFY
% def bfs(graphe, depart):
%     visites = [depart]
%     file = [depart]
%     while file:
%         sommet = file.pop(0)
%         for voisin in graphe[sommet]:
%             if voisin not in visites:
%                 visites.append(voisin)
%                 file.append(voisin)
%     return visites
% graphe = {
%     "P": ["Q", "R"], "Q": ["P", "S"], "R": ["P", "S"], "S": ["Q", "R"],
% }
% assert bfs(graphe, "P") == ["P", "Q", "R", "S"]
% def a_un_cycle(graphe, sommet, visites, parent):
%     visites.append(sommet)
%     for voisin in graphe[sommet]:
%         if voisin not in visites:
%             if a_un_cycle(graphe, voisin, visites, sommet):
%                 return True
%         elif voisin != parent:
%             return True
%     return False
% assert a_un_cycle(graphe, "P", [], None) is True
% def dfs(graphe, sommet, visites=None):
%     if visites is None:
%         visites = []
%     visites.append(sommet)
%     for voisin in graphe[sommet]:
%         if voisin not in visites:
%             dfs(graphe, voisin, visites)
%     return visites
% assert dfs(graphe, "P") == ["P", "Q", "S", "R"]
% assert dfs(graphe, "P") != bfs(graphe, "P")
% def existe_chemin(graphe, depart, arrivee):
%     return arrivee in dfs(graphe, depart)
% assert existe_chemin(graphe, "P", "S") is True
% # Un sommet isole n'est atteignable depuis aucun autre.
% graphe_isole = dict(graphe)
% graphe_isole["T"] = []
% assert existe_chemin(graphe_isole, "P", "T") is False
% END-VERIFY

\bigskip

\textbf{Exercice 2} --- \textit{Diviser pour régner} \hfill (6 points)

\begin{enumerate}
  \item (2 pts) Rappeler les trois étapes de la méthode diviser pour régner.
  \item (4 pts) Écrire une fonction récursive \lstinline{somme_dpr(liste)} qui calcule la somme des éléments d'une liste non vide par diviser pour régner (diviser en deux moitiés, additionner récursivement, puis combiner par addition).
\end{enumerate}

% BEGIN-VERIFY
% def somme_dpr(liste):
%     if len(liste) == 1:
%         return liste[0]
%     milieu = len(liste) // 2
%     return somme_dpr(liste[:milieu]) + somme_dpr(liste[milieu:])
% assert somme_dpr([1, 2, 3, 4, 5]) == 15
% assert somme_dpr([7]) == 7
% END-VERIFY

\bigskip

\textbf{Exercice 3} --- \textit{Programmation dynamique} \hfill (4 points)

On veut calculer le $n$-ieme terme de la suite de Fibonacci ($F_0=0$, $F_1=1$,
$F_{n}=F_{n-1}+F_{n-2}$).

\begin{enumerate}
  \item (1 pt) Expliquer pourquoi l'ecriture recursive naive recalcule un tres grand nombre
  de fois les memes termes.
  \item (2 pts) Ecrire une version par \textbf{programmation dynamique} qui memorise les
  resultats deja calcules, et donner $F_{30}$.
  \item (1 pt) Comparer le nombre d'appels des deux versions pour $n=20$, et dire ce que la
  memorisation echange contre ce gain.
\end{enumerate}

% BEGIN-VERIFY
% appels = {"naif": 0, "dynamique": 0}
%
% def fibonacci_naif(n):
%     appels["naif"] += 1
%     if n < 2:
%         return n
%     return fibonacci_naif(n - 1) + fibonacci_naif(n - 2)
%
% def fibonacci_dynamique(n, memo=None):
%     if memo is None:
%         memo = {}
%     appels["dynamique"] += 1
%     if n < 2:
%         return n
%     if n not in memo:
%         memo[n] = fibonacci_dynamique(n - 1, memo) + fibonacci_dynamique(n - 2, memo)
%     return memo[n]
%
% assert fibonacci_dynamique(30) == 832040
% assert fibonacci_dynamique(0) == 0 and fibonacci_dynamique(1) == 1
% appels["naif"] = appels["dynamique"] = 0
% assert fibonacci_naif(20) == fibonacci_dynamique(20) == 6765
% # Le naif fait des milliers d'appels la ou la version memorisee en fait
% # quelques dizaines.
% assert appels["naif"] > 13000
% assert appels["dynamique"] < 60
% assert appels["naif"] > 200 * appels["dynamique"]
% END-VERIFY

\bigskip

\textbf{Exercice 4} --- \textit{Recherche d'un motif : Boyer-Moore} \hfill (3 points)

\begin{enumerate}
  \item (2 pts) Expliquer le principe de l'algorithme de Boyer-Moore : dans quel sens les
  caracteres du motif sont-ils compares, et que permet la regle du \textbf{mauvais
  caractere} lorsqu'une comparaison echoue ?
  \item (1 pt) On cherche le motif \lstinline{"NSI"} dans le texte
  \lstinline{"BACNSIBAC"}. Indiquer de combien de positions le motif peut etre decale apres
  le premier echec, et pourquoi cela fait moins de comparaisons qu'une recherche naive.
\end{enumerate}

% BEGIN-VERIFY
% # Table du mauvais caractere : derniere position de chaque lettre du motif.
% def table_mauvais_caractere(motif):
%     return {lettre: indice for indice, lettre in enumerate(motif)}
%
% motif, texte = "NSI", "BACNSIBAC"
% table = table_mauvais_caractere(motif)
% assert table == {"N": 0, "S": 1, "I": 2}
% # Premiere comparaison : le motif est aligne en 0, on compare par la DROITE,
% # donc texte[2] = "C" contre motif[2] = "I". "C" n'est pas dans le motif :
% # le motif peut sauter entierement, soit un decalage de 3.
% assert texte[2] == "C" and motif[2] == "I"
% assert "C" not in table
% decalage = len(motif)
% assert decalage == 3
% # Le motif est alors aligne en 3, ou il correspond exactement.
% assert texte[3:6] == motif
% # Comparaisons faites : 1 (echec) puis 3 (succes), soit 4 -- contre 5
% # alignements testes par la recherche naive.
% naives = sum(1 for depart in range(len(texte) - len(motif) + 1))
% assert naives == 7
% assert 4 < naives
% END-VERIFY
